\documentclass[12pt,a4paper]{article}
\usepackage[UTF8]{ctex}
\usepackage{geometry}
\usepackage{amsmath,amssymb}
\usepackage{booktabs,array,tabularx}
\usepackage{enumitem,hyperref,fancyhdr,xcolor,setspace}
\usepackage{ragged2e}

\geometry{left=2.5cm,right=2.5cm,top=2.5cm,bottom=2.5cm}
\setlength{\headheight}{15pt}
\setstretch{1.5}
\setlength{\emergencystretch}{3em}
\hfuzz=20pt
\sloppy
\hypersetup{hidelinks}

\pagestyle{fancy}
\fancyhf{}
\fancyhead[C]{说明书摘要}
\fancyfoot[C]{\thepage}
\renewcommand{\headrulewidth}{0pt}

\begin{document}

\begin{center}
    {\Huge\heiti 说明书摘要} \\[0.8em]
\end{center}

\noindent
本发明公开了一种适用于长距离量子中继器的分布式纠缠纯化与纠缠交换的方法、节点装置及系统。该方法在各中继节点维持锚定纠缠对与候选牺牲纠缠对，利用基于保真度估计、纠缠对年龄与链路重建代价的综合评分触发滚动式小批量纯化；并将纯化窗口与纠缠交换窗口重叠为流水线，以在等待经典确认的间隙继续生成候选纠缠对，从而提高吞吐。通过在经典信道上传输包含纠缠对标识、层级、会话与租约的压缩综合征消息，实现异步链路下的局部锁定、冲突避免与局部失败恢复，使得在少量物理量子比特和有限量子存储相干时间条件下能够持续输出高保真长距离纠缠对。

\vspace{0.6em}
\noindent
摘要附图：图1为分布式纠缠纯化架构及纯化/交换窗口重叠时序示意图。

\end{document}

